\section{A complete polynomial test on certified treewidth-two projections}
\label{sec:treewidth-two}

\subsection{The structural restriction that makes determinant one sufficient}

The general determinant-one implication is false: the Conway and
Kinoshita--Terasaka knots both survive that scalar test.  Here we add an
explicit graph hypothesis to make it a complete decision rule.  Let $G_D$
be the projection graph of a validated classical one-component diagram $D$:
vertices are crossings and edges are the arcs between crossings.  This is
\emph{not} a Tait graph or a knot-complement triangulation graph.
Loops and multiple edges do not affect the presence of a $K_4$ minor.

The structural theorem of Ru\'e, Thilikos and Velona
\cite{rue-thilikos-velona} classifies links admitting a $K_4$-minor-free
projection as disjoint and connected sums of two-strand torus links.
For a knot, this reduces to connected sums of $T(2,q_i)$ with odd $q_i$,
with trivial factors allowed.  Indeed a connected sum of links with
$c_1,c_2$ components has $c_1+c_2-1$ components, so a two-component torus
link cannot occur in a decomposition resulting in one component.

\begin{proposition}
For a classical knot diagram $D$ with $\operatorname{tw}(G_D)\le2$,
\[
 D\text{ represents the unknot}\quad\Longleftrightarrow\quad
 \det(D)=1.
\]
\end{proposition}
\begin{proof}
Treewidth at most two is equivalent to excluding a $K_4$ minor.  Apply the
classification above.  For odd $q$, the Alexander polynomial of $T(2,q)$,
up to a Laurent unit and mirroring, is
$(t^{|q|}+1)/(t+1)$, whose value at $t=-1$ has absolute value $|q|$.
The Alexander polynomial multiplies under connected sum, hence
$\det(D)=\prod_i |q_i|$.  This product is one exactly when all factors
are trivial.  Conversely every unknot has determinant one.
\end{proof}

This uses the given projection, with no supplied braid or Montesinos
presentation.  A rejected projection does not imply that its knot type
has no other diagram in the class.  In particular, it is not a
knottedness certificate.

Bodlaender, Burton, Fomin and Grigoriev \cite{bodlaender-treewidth-two}
give a stronger, linear-time recognition algorithm on the same class,
using generalized diagrams with twist-labelled double edges.  Our
implementation uses the classification and an exact determinant instead;
it does not implement their linear-time algorithm or inherit its time
bound.  This separates a directly checkable implementation from a future
optimization of its arithmetic.

\subsection{A graph certificate before any positive invariant verdict}

The producer in \code{fastunknot/treewidth\_two.py} builds the simple
projection adjacency sets.  It queues vertices of degree at most two.
Deleting a degree-two vertex joins its two remaining neighbours; existing
parallel edges coalesce.  It records the deletion order.  Every operation
preserves the property of having treewidth at most two in both directions.
For degree two, suppression is an edge contraction followed by removal of
redundant parallel edges, so it cannot increase treewidth.  Conversely, a
width-two decomposition of the remaining graph has a bag containing the
new neighbour edge; attach the deleted vertex in a new bag of size three.
The cases of degree zero or one are immediate.

Every nonempty graph of treewidth at most two has a vertex of degree at
most two.  Consequently elimination either deletes everything and supplies
a certificate, or stops at a graph of minimum degree at least three and
rejects the input projection.  Deletion never increases a remaining
vertex's degree: a possible fill edge replaces the incident edge just
removed.  This justifies the queue implementation.  Each deletion examines
at most two neighbours and creates at most one fill edge.  The graph phase
uses $O(n)$ adjacency operations and $O(n)$ words; the recorded order uses
$O(n\log(n+1))$ bits.

A verifier does not trust the producer's queue or adjacency state.  It
reconstructs projection edges from PD labels, checks that the order is a
permutation of all crossings, and replays deletion and fill on an edge set.
It rejects any step with more than two remaining neighbours.  The simple
independent replay is polynomial rather than linear.  Crucially, it must
succeed \emph{before} an exact determinant of one is allowed to certify an
unknot.  Certificates also bind the crossing count, schema version, method
and exact verdict.  The width witness is a sufficient structural certificate,
not an implementation of a formal proof of the imported classification.

\subsection{Exact integer arithmetic and independent verification}

The producer builds the signed checkerboard graph using the existing
PD corner convention.  For either shading, the absolute signed
spanning-tree sum is the knot determinant.  We choose the shading with
fewer vertices, discard loops, sum parallel signed weights and use the
articulation-block product from Section~\ref{sec:tait-blocks}.
The reduced signed Laplacian of every nontrivial block is evaluated by
fraction-free Bareiss elimination with exact division and row exchanges
when a leading pivot is zero.  Bridge blocks contribute just their integer
weight.  The shared determinant routine was moved to
\code{integer\_determinant.py}; the marked-continuation backend retains its
existing interface and algorithm.

No modular residue is promoted to an exact determinant.  A cancelled
parallel sum and a zero leading pivot require exact handling even though a
validated classical knot ultimately has an odd, nonzero determinant.
The stored determinant is hexadecimal, so large composite knots do not
hit Python's decimal-conversion ceiling when their certificates are
serialized as JSON.

The independent certificate verifier avoids the checkerboard construction
and Bareiss entirely.  It constructs Wirtinger arcs by joining only the
overpassing arc labels at each crossing.  At $t=-1$, every Fox row is
\[
 2x_{\rm over}-x_{\rm under,1}-x_{\rm under,2},
\]
independently of crossing sign and orientation.  It removes one row and
one column and evaluates the resulting cofactor using exact rational
Gaussian elimination.  Its absolute value must equal the reported
integer, and the verdict must be precisely the determinant-one decision.
This provides a separate graph representation, knot-matrix construction
and elimination algorithm for replay.

With $n$ crossings, checkerboard matrices have $O(n)$ rows and integer
entries of size $O(n)$.  Hadamard bounds give $O(n\log(n+1))$ bits for
all relevant minors, hence for Bareiss entries; numerator products and
exact divisors remain polynomial in bit length.  Summing cubic arithmetic
cost over blocks gives at most $O(n^3)$ integer arithmetic operations.
Thus the uncapped certificate producer has a polynomial bit bound on this
class, with a conservative dense memory bound also polynomial.  This is
weaker than the cited linear algorithm, but independent of the size of any
Khovanov complex.

\subsection{Adaptive integration and the exact scope of the guarantee}

The CLI option \code{--treewidth-two}, or Python option
\code{use\_treewidth\_two=True}, enables this stage after checked source
certificates and before the usual Seifert and simplification stages.
The local allowance defaults to $0.1$ seconds and includes lazy module
loading, graph certification and determinant work.  Its callback first
checks the global deadline, then the local one; graph loops and determinant
rows poll cooperatively.  A local timeout or allocation failure records a
skipped probe and continues the existing recognizer.  A projection outside
the class also continues.  A global timeout retains \code{UNKNOWN}.
No interrupted order or partially multiplied determinant is published as
a verdict.  The local interval and discarded work remain part of the
complete query's global elapsed time.

The standalone producer has no local cap.  An uncapped local query in the
Python recognizer sets \code{treewidth\_\allowbreak two\_\allowbreak seconds}
to \code{None}.
The polynomial claim belongs to the uncapped
restricted-class procedure.  A capped probe that falls back to Khovanov
does not give the whole pipeline a polynomial guarantee, even on inputs
whose projection lies in the class.  The optional stage also does not
change the general exponential upper bound or establish a general
subexponential recognizer.

\subsection{Validation and whole-query measurements}

Eight added tests cover: an independent branch-set search for $K_4$ minors
on all 64 four-vertex graphs and 25 larger graphs; all 170 signed
one-component two-strand braids through seven crossings against full
cube-of-resolutions recognition; random braid and named-diagram determinants
against exact Alexander polynomials; mirrored and relabelled diagrams;
forged certificates; CLI dispatch; early and late cancellation; and local
fallback versus global exhaustion.  The Conway, Kinoshita--Terasaka and
stored determinant-one grid knots explicitly test that the graph restriction
cannot be omitted.  A connected sum of 10,000 trefoils produces determinant
$3^{10000}$ and a JSON-safe certificate on 30,000 crossings.  The full suite
result is recorded in \code{data/treewidth-two-integrated-tests.txt}.

\code{benchmark\_treewidth\_two.py} measures validation plus full recognition
on fresh, source-free PD inputs.  Fourteen inputs each receive one excluded
warm-up and seven shuffled paired rounds: default, an identical default
control, and the optional stage.  The interpreter is warm; imports are not
part of the measured rounds.  All 294 measured queries finish with matching
verdicts.  Median paired speedups on mixed-sign two-strand unknot diagrams
of 101, 501 and 2001 crossings are respectively $1.085$, $1.142$ and $1.126$.
The default/control ratios range from $0.954$ to $1.038$.

Other inputs show why the stage remains optional.  The trefoil has a
$0.606$ speedup ratio, tree medials range from $0.756$ to $0.838$, and the
two 120-crossing trefoil sums have ratios near $0.67$: existing Seifert
certificates already decide these cheaply.  Outside-class inputs range
from $0.862$ to $1.018$.  These measurements support a modest improvement
on a specific family and a new independently certified input-class bound,
not a default-wide performance improvement.  Exact inputs, execution order,
individual timings, methods and certificates are retained in
\code{results/treewidth\_two\_20261008.json}.
